\pdfoutput=1
\documentclass[10pt,twocolumn]{article}

\usepackage[utf8]{inputenc}
\usepackage[T1]{fontenc}
\usepackage{times}
\usepackage{microtype}
\usepackage{amsmath,amssymb,amsfonts,amsthm}
\usepackage{graphicx}
\usepackage{booktabs}
\usepackage{multirow}
\usepackage{array}
\usepackage{enumitem}
\usepackage{xcolor}
\usepackage{caption}
\usepackage{subcaption}
\usepackage{hyperref}
\usepackage{geometry}
\usepackage{titlesec}

\geometry{
  paper=a4paper,
  left=16mm,
  right=16mm,
  top=19mm,
  bottom=23mm,
  columnsep=7mm
}

\hypersetup{
  colorlinks=true,
  linkcolor=blue,
  citecolor=blue,
  urlcolor=blue
}

\titleformat{\section}{\large\bfseries}{\thesection.}{0.5em}{}
\titleformat{\subsection}{\normalsize\bfseries}{\thesubsection.}{0.5em}{}
\titleformat{\subsubsection}{\normalsize\itshape}{\thesubsubsection.}{0.5em}{}

\setlength{\parindent}{1em}
\setlength{\parskip}{0pt}
\setlist[itemize]{leftmargin=1.15em,itemsep=1pt,topsep=2pt}
\setlist[enumerate]{leftmargin=1.25em,itemsep=1pt,topsep=2pt}

\newtheorem{proposition}{Proposition}
\newcommand{\method}{SemEq}
\newcommand{\ie}{i.e.,}
\newcommand{\eg}{e.g.,}
\newcommand{\etal}{et al.}
\newcommand{\pending}{--}
\newcommand{\placeholderfigure}[1]{
  \begin{center}
  \fbox{
    \begin{minipage}[c][0.18\textheight][c]{0.92\linewidth}
    \centering #1
    \end{minipage}
  }
  \end{center}
}

\begin{document}

\twocolumn[
\begin{center}
{\LARGE \bf \method: Equilibrium-Guided Semantic Coordination for\\
Training-Free High-Resolution Video Generation}

\vspace{0.6em}
{\large Anonymous ECCV Submission}

\vspace{0.8em}
\end{center}

\begin{abstract}
Training-free high-resolution video generation requires a pretrained video diffusion model to operate far beyond the spatial token distribution observed during pretraining. The central challenge is not only computational. A high-resolution denoising step must preserve prompt-level structure while recovering native-resolution visual evidence such as facial details, object boundaries, and fine textures. Existing inference-time rules usually resolve this conflict with fixed attention choices or fixed fusion coefficients, which cannot adapt to prompt content, denoising stage, token location, or cache reliability. This paper introduces \method, an equilibrium-guided semantic coordination principle for high-resolution video diffusion. At each cross-attention layer, \method observes two forms of semantic evidence: contextual evidence for scene-level structure and native-grid evidence for high-frequency detail. Their reliability is measured from attention alignment and response confidence, and the final semantic authority is obtained by solving a Nash bargaining objective under a shared authority constraint. The resulting closed-form rule is adaptive, bounded, interpretable, and requires no training, no additional data, and no model-weight modification. We further provide a matched-seed evaluation protocol and ablations to verify whether equilibrium coordination improves 1080P quality, consistency, and runtime compared with fixed inference strategies.
\end{abstract}

\vspace{1.1em}
]

\section{Introduction}

Text-to-video diffusion models have made rapid progress in synthesizing coherent short videos from natural language prompts. However, their learned generation prior is usually tied to a moderate native resolution. When the same denoiser is directly applied to 1080P synthesis, the number of spatial tokens increases dramatically. This shift often causes blurred details, unstable small objects, distorted faces, repeated structures, and weak temporal consistency. The problem is especially visible in close-up shots, dense backgrounds, and videos that require both large-scale layout and local texture fidelity.

A natural strategy is to use inference-time refinement without retraining the model. High-resolution inference can reuse a pretrained generator and avoid costly high-resolution video datasets. Yet the inference rule must decide how semantic information is injected into the high-resolution latent. Contextual attention preserves prompt-level meaning, identity, object count, and scene layout. Native-grid evidence preserves local visual fidelity under a spatial regime closer to the model's training distribution. These two forms of evidence are complementary, but they are not equally reliable at every token or denoising step. A fixed rule may over-trust contextual semantics and wash out details, or over-trust native-grid evidence and fragment the scene.

This paper argues that high-resolution refinement should be formulated as a semantic coordination problem. At each cross-attention layer, the denoiser must assign one unit of semantic authority to multiple forms of evidence. This authority determines how strongly each evidence form influences the final hidden representation. Instead of setting this authority by hand, we seek a principle that is adaptive, training-free, and stable during diffusion sampling.

We propose \method, an equilibrium-guided semantic coordination method. The key idea is to treat contextual semantic evidence and native-grid visual evidence as two cooperative terms sharing a constrained semantic resource. Their current reliability is measured from attention behavior: alignment indicates whether an evidence response is compatible with the requesting token, while confidence indicates whether the response is active relative to the layer context. Given these reliability values, \method solves a Nash bargaining objective and obtains a closed-form authority allocation. The solution increases the authority of more reliable evidence while avoiding collapse to a single term when both remain useful.

This formulation is different from a heuristic weighted sum. The mixture coefficient is not a manually selected constant. It is the equilibrium solution of a defined semantic coordination objective. The method therefore provides a clear rationale for when contextual structure should dominate and when native-grid detail should be trusted. It is also compatible with cached inference: expensive contextual evidence can be reused across denoising steps, while the equilibrium rule still evaluates whether the cached response deserves authority at the current step.

\method is designed for practical deployment. It does not finetune the video model, add a learned controller, or require extra annotations. The added computation consists of lightweight reliability observation and a closed-form allocation inside attention layers. This makes the approach suitable for 1080P generation where both quality and runtime matter.

Our contributions are:
\begin{itemize}
    \item We formulate training-free high-resolution video refinement as equilibrium-guided semantic coordination rather than fixed attention selection.
    \item We introduce attention-grounded reliability measures based on query-response alignment and response confidence, requiring no learned scoring network.
    \item We derive a closed-form Nash semantic authority allocation and stabilize it with a bounded mapping suitable for diffusion inference.
    \item We provide an experimental protocol, ablations, and visualization design for evaluating 1080P quality, temporal consistency, and runtime under matched seeds.
\end{itemize}

\section{Related Work}

\subsection{Video Diffusion Models}

Diffusion models have become a dominant approach for visual generation \cite{ho2020ddpm,rombach2022ldm}. Transformer-based denoisers further improve scalability by using attention to model long-range dependencies \cite{vaswani2017attention,peebles2023dit}. In video generation, attention is crucial because appearance, motion, and prompt semantics must remain coordinated across both space and time. However, attention cost grows quickly with token count, and high-resolution video inference introduces a severe token-distribution shift.

\subsection{Training-Free High-Resolution Inference}

Training-free high-resolution inference aims to improve output resolution without collecting high-resolution data or finetuning large generative models. Existing strategies include latent upsampling, noise-based refinement, local attention, token reuse, and cached computation \cite{meng2021sdedit,bolya2023tome,ma2024deepcache}. These methods are attractive because they can be applied to pretrained models. Their limitation is that all decisions must be made during inference, making the rule that coordinates contextual semantics and native-grid evidence especially important.

\subsection{Locality and Semantic Consistency}

Local computation is useful because small neighborhoods are closer to the model's native training regime and better preserve high-frequency detail \cite{liu2021swin,hassani2023nat}. Nevertheless, video generation cannot be solved by local evidence alone. A prompt specifies identity, count, layout, background, attributes, and motion. These constraints require contextual semantic coordination. The central challenge is therefore not whether to use contextual or native-grid evidence, but how to allocate authority between them adaptively.

\subsection{Caching in Diffusion Sampling}

Caching accelerates diffusion sampling by reusing intermediate computations across steps \cite{ma2024deepcache}. In high-resolution video generation, contextual evidence can be expensive, so reusing it is valuable. However, cached responses may become stale as the latent evolves. \method treats cached contextual evidence as an evidence term whose authority must still be justified by current reliability. Thus, caching improves efficiency while equilibrium coordination controls whether the cached response remains trustworthy.

\subsection{Nash Equilibrium and Bargaining}

Nash bargaining provides an axiomatic solution for allocating a shared resource among cooperative terms with positive utilities \cite{nash1950bargaining,nash1951noncooperative}. Maximizing the Nash product balances efficiency and non-collapse: a term with higher utility receives more resource, but the solution discourages assigning all authority to one term when multiple terms remain useful. This property matches high-resolution semantic coordination, where both contextual structure and native-grid detail are needed for stable 1080P video synthesis.

\section{Method}

\subsection{High-Resolution Refinement}

Given a text prompt $c$, a pretrained video diffusion model generates or refines a latent video. Let $z_t$ denote the latent at denoising step $t$. In a cross-attention layer, visual queries interact with textual keys and values to inject prompt semantics. At high resolution, a visual token requires both contextual semantic structure and native-grid visual detail.

For token $i$, let
\begin{equation}
    O_c(i), O_n(i) \in \mathbb{R}^{H \times D}
\end{equation}
be two candidate cross-attention responses before output projection. The subscript $c$ denotes contextual semantic evidence, and $n$ denotes native-grid evidence. $H$ is the number of attention heads and $D$ is the head dimension. The final response is
\begin{equation}
    O(i)=a(i)O_c(i)+(1-a(i))O_n(i),
    \label{eq:mix}
\end{equation}
where $a(i)$ is the semantic authority assigned to contextual evidence. The goal is to determine $a(i)$ adaptively and without training.

\subsection{Semantic Evidence}

\noindent\textbf{Contextual semantic evidence.}
This evidence preserves prompt-level structure, object identity, object count, long-range layout, and broad temporal consistency. It can be obtained from full-context cross-attention or from cached full-context attention responses. It is especially important during early denoising and in regions where global layout or identity constraints dominate.

\noindent\textbf{Native-grid visual evidence.}
This evidence preserves local texture, boundaries, and high-frequency details under a spatial regime closer to the model's native training distribution. It can be obtained from native-scale windowed attention or another local evidence construction. It is especially important during later denoising, close-up synthesis, and regions with fine structures.

Both evidence forms are produced by the same pretrained model at inference time. The proposed method does not create separate networks. It only decides how much authority each evidence form should receive.

\subsection{Attention-Grounded Reliability}

The equilibrium rule requires positive reliability values. \method computes them from attention responses directly. For evidence type $x\in\{c,n\}$, let $Q_x(i)$ be the corresponding query tensor and $O_x(i)$ be its candidate response.

First, query-response alignment measures whether the response direction is compatible with the requesting token:
\begin{equation}
    A_x(i)=
    \frac{1+\cos(\operatorname{vec}(O_x(i)),
    \operatorname{vec}(Q_x(i)))}{2}.
    \label{eq:alignment}
\end{equation}
This term lies in $[0,1]$. A high value indicates that the semantic response is directionally consistent with the token representation.

Second, response confidence measures whether the response is active relative to the layer context:
\begin{equation}
    C_x(i)=
    \sigma\left(
    \frac{\operatorname{RMS}(O_x(i))}
    {\operatorname{mean}_{j}\operatorname{RMS}(O_x(j))+\epsilon}
    -1
    \right).
    \label{eq:confidence}
\end{equation}
This term prevents weak responses from receiving high authority only because they are directionally aligned.

The reliability value is
\begin{equation}
    r_x(i)=\pi_x\left(A_x(i)+C_x(i)+\epsilon\right),
    \label{eq:reliability}
\end{equation}
where $\pi_x$ is a prior. By default, $\pi_c=1.25$ and $\pi_n=1.0$. The contextual prior protects prompt-level structure, while native-grid evidence can still dominate when its observed reliability is stronger.

\subsection{Nash Semantic Coordination}

For each token, contextual and native-grid evidence share one unit of semantic authority. The feasible set is
\begin{equation}
    \Delta=\{(a_c,a_n)\mid a_c>0,\; a_n>0,\; a_c+a_n=1\}.
\end{equation}
Let $a_c=a$ and $a_n=1-a$. We define the Nash semantic coordination objective as
\begin{equation}
    \max_{a\in(0,1)}
    r_c(i)\log(a+\epsilon)
    +
    r_n(i)\log(1-a+\epsilon).
    \label{eq:nash}
\end{equation}
This objective maximizes the log Nash product under the shared-authority constraint.

\begin{proposition}
Ignoring the small numerical constant $\epsilon$, Eq.~\eqref{eq:nash} has the unique optimum
\begin{equation}
    a^*(i)=
    \frac{r_c(i)}
    {r_c(i)+r_n(i)}.
    \label{eq:closed}
\end{equation}
\end{proposition}

\noindent\textbf{Proof.}
The objective is strictly concave because
\begin{equation}
    \frac{\partial^2}{\partial a^2}
    \left[
    r_c\log a+r_n\log(1-a)
    \right]
    =
    -\frac{r_c}{a^2}
    -\frac{r_n}{(1-a)^2}<0.
\end{equation}
The first-order condition is
\begin{equation}
    \frac{r_c}{a}
    -
    \frac{r_n}{1-a}=0,
\end{equation}
which gives Eq.~\eqref{eq:closed}. Strict concavity makes the solution unique. \hfill $\square$

The solution is interpretable. If contextual reliability increases, contextual authority increases. If native-grid reliability increases, native-grid authority increases. Equal reliability values produce equal unconstrained authority. Therefore, the coefficient in Eq.~\eqref{eq:mix} is not a heuristic constant; it is the equilibrium solution of a defined semantic coordination objective.

\subsection{Bounded Authority Mapping}

Unbounded authority may be unstable during diffusion sampling. We therefore temper and bound the equilibrium ratio. First, reliability values are temperature-scaled:
\begin{equation}
    \tilde{r}_x(i)=r_x(i)^{1/\tau}.
\end{equation}
Then the equilibrium ratio is
\begin{equation}
    q(i)=
    \frac{\tilde{r}_c(i)}
    {\tilde{r}_c(i)+\tilde{r}_n(i)+\epsilon}.
    \label{eq:ratio}
\end{equation}
Finally, the authority is mapped to a stable interval:
\begin{equation}
    a(i)=
    a_{\min}
    +(a_{\max}-a_{\min})q(i).
    \label{eq:bounded}
\end{equation}
The default setting is $a_{\min}=0.65$ and $a_{\max}=0.98$. The lower bound protects prompt-level structure, while the upper bound preserves a minimum native-grid contribution.

\subsection{Cache-Compatible Inference}

Contextual evidence can be expensive at high resolution. \method supports cached inference by refreshing contextual evidence every $P$ steps and reusing it between refreshes. Let $t'$ be the most recent refresh step. At step $t$, the response is
\begin{equation}
    O_t(i)=
    a_t(i)O_{c,t'}(i)
    +(1-a_t(i))O_{n,t}(i).
    \label{eq:cached}
\end{equation}
Native-grid evidence remains current, while contextual evidence can be cached. Crucially, cached evidence does not receive fixed authority. Its authority is still determined by Eq.~\eqref{eq:reliability}--Eq.~\eqref{eq:bounded}, so stale contextual responses can be weakened when current native-grid evidence becomes more reliable.

\begin{figure*}[t]
\placeholderfigure{Replace this box with the final method figure. Recommended layout: prompt and high-resolution latent, contextual semantic evidence, native-grid visual evidence, reliability observation, Nash semantic coordination, final 1080P video.}
\caption{Overview of \method. The method allocates semantic authority between contextual structure and native-grid detail through a Nash equilibrium rule.}
\label{fig:overview}
\end{figure*}

\subsection{Algorithm}

\begin{enumerate}
    \item Generate or initialize a video latent from prompt $c$.
    \item Upsample or refine the latent to the target 1080P grid.
    \item At each cross-attention layer, compute contextual response $O_c$ and native-grid response $O_n$.
    \item Compute reliability values $r_c$ and $r_n$ from alignment and confidence.
    \item Solve the closed-form Nash authority allocation.
    \item Fuse the two responses using Eq.~\eqref{eq:mix}.
    \item Continue denoising and decode the final high-resolution video.
\end{enumerate}

\section{Experiments}

\subsection{Evaluation Protocol}

The main goal is to test whether equilibrium-guided semantic coordination improves high-resolution generation compared with fixed inference rules. All methods should use the same pretrained model, prompt set, random seeds, precision, frame number, denoising steps, guidance scale, and target resolution. When possible, the same low-resolution initialization should be reused across methods.

We evaluate at $1920\times1088$, reported as the 1080P setting because the height is rounded to satisfy latent-grid divisibility. Unless otherwise stated, videos contain 81 frames, high-resolution refinement uses 35 denoising steps, and the cache interval is $P=2$.

\subsection{Metrics}

Automatic evaluation follows VBench-style dimensions \cite{huang2024vbench}: subject consistency, background consistency, motion smoothness, dynamic degree, aesthetic quality, imaging quality, overall consistency, and runtime. The main table should report the average over a fixed prompt set. A single generated video is useful for sanity checking but should not be used as the final evidence for superiority.

\subsection{Main Results}

\begin{table*}[t]
\centering
\caption{Main 1080P comparison. Fill the numerical entries after running the same prompt set and matched-seed protocol for all methods. Higher is better except for runtime.}
\label{tab:main}
\begin{tabular}{lcccccccc}
\toprule
Method & Subject & Background & Motion & Dynamic & Aesthetic & Imaging & Overall & Time \\
\midrule
Direct high-resolution inference & \pending & \pending & \pending & \pending & \pending & \pending & \pending & \pending \\
Native-grid evidence only & \pending & \pending & \pending & \pending & \pending & \pending & \pending & \pending \\
Contextual evidence only & \pending & \pending & \pending & \pending & \pending & \pending & \pending & \pending \\
Fixed semantic coordination & \pending & \pending & \pending & \pending & \pending & \pending & \pending & \pending \\
\method\ without cache & \pending & \pending & \pending & \pending & \pending & \pending & \pending & \pending \\
\method\ with cache, $P=2$ & \pending & \pending & \pending & \pending & \pending & \pending & \pending & \pending \\
\bottomrule
\end{tabular}
\end{table*}

The most important comparison is between fixed semantic coordination and full \method. The fixed variant uses the same contextual and native-grid evidence but removes the Nash equilibrium allocation. Therefore, an improvement in this comparison directly supports the proposed coordination principle.

\subsection{Ablation Study}

\begin{table}[t]
\centering
\caption{Ablation on equilibrium-guided semantic coordination.}
\label{tab:ablation}
\begin{tabular}{lccc}
\toprule
Variant & Imaging & Consistency & Overall \\
\midrule
Fixed coordination & \pending & \pending & \pending \\
Alignment only & \pending & \pending & \pending \\
Confidence only & \pending & \pending & \pending \\
No contextual prior & \pending & \pending & \pending \\
Unbounded authority & \pending & \pending & \pending \\
Full \method & \pending & \pending & \pending \\
\bottomrule
\end{tabular}
\end{table}

This ablation tests whether the equilibrium formulation is necessary. Removing the Nash rule tests the value of principled semantic coordination. Removing alignment or confidence tests the reliability design. Removing the contextual prior tests whether protecting prompt-level structure is beneficial. Removing bounds tests whether the unconstrained equilibrium is stable enough for diffusion inference.

\subsection{Cache Interval Analysis}

\begin{table}[t]
\centering
\caption{Cache interval analysis. Runtime should be measured on the same GPU and with the same precision.}
\label{tab:cache}
\begin{tabular}{lccc}
\toprule
Setting & Overall & Runtime & Speedup \\
\midrule
No cache & \pending & \pending & 1.0$\times$ \\
$P=2$ & \pending & \pending & \pending \\
$P=5$ & \pending & \pending & \pending \\
$P=8$ & \pending & \pending & \pending \\
\bottomrule
\end{tabular}
\end{table}

Larger cache intervals reduce computation but may make contextual evidence stale. \method is expected to preserve quality under moderate caching because cached responses must still pass the reliability-based authority allocation.

\subsection{Qualitative Comparison}

\begin{figure*}[t]
\placeholderfigure{Replace this box with full-frame comparisons and zoomed crops. Recommended cases: face close-up, dense background, moving object boundary, small texture, and long-range layout.}
\caption{Qualitative comparison template. The final figure should show both full-frame semantic structure and zoomed local evidence.}
\label{fig:qualitative}
\end{figure*}

Qualitative examples should emphasize the failure modes that the method is designed to address. Native-grid-heavy inference may produce sharp local details but unstable global structure. Context-heavy inference may preserve layout but smooth out local evidence. Fixed coordination may work for some prompts but fail when the required balance changes across regions or denoising stages. \method should demonstrate stronger balance because semantic authority is recomputed from current attention evidence.

\section{Discussion}

\subsection{Why This Is Not Merely a Weighted Sum}

At the implementation level, the final cross-attention response is a mixture of two tensors. However, the contribution of this paper is not the existence of a mixture. The contribution is the rule that determines the mixture coefficient. A fixed coefficient is an engineering choice. \method instead defines semantic evidence, reliability values, a feasible authority set, and a Nash objective. The coefficient is the solution to this objective. This gives the method a rationality principle and makes it testable through ablation.

\subsection{Limitations}

\method depends on the quality of its reliability observation. Alignment and confidence are efficient and available inside attention layers, but they may not capture all forms of semantic reliability. Future work could include temporal reliability, motion-aware evidence, prompt-token saliency, or uncertainty estimation. The method also assumes that both evidence forms are meaningful. When the pretrained model lacks the ability to represent a concept, equilibrium coordination cannot recover missing knowledge.

\section{Conclusion}

We presented \method, an equilibrium-guided semantic coordination method for training-free high-resolution video generation. The method treats contextual structure and native-grid detail as complementary evidence forms sharing one unit of semantic authority. Their reliability is computed directly from attention behavior, and the final authority is obtained through a closed-form Nash bargaining solution. This provides a principled alternative to fixed inference rules, remains compatible with cached 1080P refinement, and offers a clear experimental path for validating equilibrium coordination in high-resolution video synthesis.

\begin{thebibliography}{99}

\bibitem{vaswani2017attention}
A. Vaswani, N. Shazeer, N. Parmar, J. Uszkoreit, L. Jones, A. N. Gomez, L. Kaiser, and I. Polosukhin.
Attention is all you need.
In \emph{NeurIPS}, 2017.

\bibitem{ho2020ddpm}
J. Ho, A. Jain, and P. Abbeel.
Denoising diffusion probabilistic models.
In \emph{NeurIPS}, 2020.

\bibitem{rombach2022ldm}
R. Rombach, A. Blattmann, D. Lorenz, P. Esser, and B. Ommer.
High-resolution image synthesis with latent diffusion models.
In \emph{CVPR}, 2022.

\bibitem{peebles2023dit}
W. Peebles and S. Xie.
Scalable diffusion models with Transformers.
In \emph{ICCV}, 2023.

\bibitem{meng2021sdedit}
C. Meng, Y. He, Y. Song, J. Song, J. Wu, J. Zhu, and S. Ermon.
SDEdit: Guided image synthesis and editing with stochastic differential equations.
In \emph{ICLR}, 2022.

\bibitem{liu2021swin}
Z. Liu, Y. Lin, Y. Cao, H. Hu, Y. Wei, Z. Zhang, S. Lin, and B. Guo.
Swin Transformer: Hierarchical vision Transformer using shifted windows.
In \emph{ICCV}, 2021.

\bibitem{hassani2023nat}
A. Hassani and H. Shi.
Neighborhood attention Transformer.
In \emph{CVPR}, 2023.

\bibitem{huang2024vbench}
Z. Huang, Y. He, J. Yu, F. Zhang, C. Si, Y. Jiang, Y. Zhang, T. Wu, Q. Jin, N. Chanpaisit, Y. Wang, X. Chen, L. Wang, D. Lin, Y. Qiao, and Z. Liu.
VBench: Comprehensive benchmark suite for video generative models.
In \emph{CVPR}, 2024.

\bibitem{nash1950bargaining}
J. Nash.
The bargaining problem.
\emph{Econometrica}, 18(2):155--162, 1950.

\bibitem{nash1951noncooperative}
J. Nash.
Non-cooperative games.
\emph{Annals of Mathematics}, 54(2):286--295, 1951.

\bibitem{bolya2023tome}
D. Bolya and J. Hoffman.
Token merging for fast stable diffusion.
In \emph{CVPR Workshop}, 2023.

\bibitem{ma2024deepcache}
X. Ma, G. Fang, and X. Wang.
DeepCache: Accelerating diffusion models for free.
In \emph{CVPR}, 2024.

\end{thebibliography}

\end{document}
